\chapter*{Code Appendix}
\addcontentsline{toc}{chapter}{Code Appendix}
\label{ch:code_appendix}

This appendix contains the source code listings referenced throughout the implementation chapter. The code is organized by component and includes detailed comments explaining key functionality.

\usubsection{Dataset Processing Script}
\begin{minted}[%
    framesep=2mm,
    baselinestretch=1.2,
    bgcolor=LightGray,
    fontsize=\footnotesize,
    breaklines
    ]{python}
    import csv
    import sys
    import json
    import os

    # Creates a new csv file named 'edited_csv' that doesn't contain empty rows, or rows with empty fields
    def csv_cleanup(csvfilename):
        with open(csvfilename, newline='') as csvfile:
            with open('edited_csv', 'w', newline='') as edited_csvfile:
                original = csv.reader(csvfile, delimiter=';')
                edited = csv.writer(edited_csvfile, delimiter=';')
                # fieldnames = next(original)
                # fieldnames.insert(0, "#")
                # edited.writerow(fieldnames)
                # i = 1
                for row in original:
                    if row and any(row) and any(field.strip() for field in row):
                        # row.insert(0, i)
                        edited.writerow(row)
                        # i += 1

    # Creates a json file from the edited csv file with an incremental integer as keys
    def csv_to_json(filename):
        data_dict = {}
        with open('edited_csv', newline='') as csvfile:
            edited = csv.DictReader(csvfile, delimiter=';')
            key = 1
            for row in edited: 
                data_dict[key] = row
                key += 1
        with open(filename+'.json', 'w') as jsonfile:
            jsonfile.write(json.dumps(data_dict, indent = 4))



    csvfilename = sys.argv[1]
    filename = csvfilename.split('/')[-1].split('.')[0]

    csv_cleanup(csvfilename)
    csv_to_json(filename)
    os.remove('edited_csv')
\end{minted}

\usubsection{Sensor Simulator main application}
\begin{minted}[%
    framesep=2mm,
    baselinestretch=1.2,
    bgcolor=LightGray,
    fontsize=\footnotesize,
    breaklines,
    label=lst:sensor simulator main application
    ]{python}
    import socket
    import time
    import subprocess
    import paho.mqtt.client as mqtt_client

    port = 1883  # Default port for MQTT communication
    topic = "collect-data"  # Topic to subscribe to for data collection
    hostname = socket.gethostname()  # Get the hostname of the machine running this script. Machine (container) hostname is enforced later-on by docker compose 
    client_id = 'subscribe-{}'.format(hostname)  # Unique client ID for MQTT connection
    # broker = 'localhost'  # Uncomment this line for local broker testing
    broker = 'host.docker.internal'  # Use when running in a Docker environment

    def connect_mqtt():
        """
        Connects to the MQTT broker and sets up event handlers for connect, disconnect, and message events.
        """

        def on_connect(client, userdata, flags, rc):
            """
            Callback triggered upon connecting to the MQTT broker.
            """
            if rc == 0:
                print("Connected to MQTT Broker!")
                client.subscribe(topic)  # Subscribe to the specified topic
            else:
                print("Failed to connect, return code %d\n", rc)
        
        def on_disconnect(client, userdata, rc):
            """
            Callback triggered when the MQTT client disconnects from the broker.
            Implements an exponential backoff strategy for reconnection attempts.
            """
            FIRST_RECONNECT_DELAY = 1
            RECONNECT_RATE = 2
            MAX_RECONNECT_COUNT = 12
            MAX_RECONNECT_DELAY = 60

            print("Disconnected with result code: %s", rc)
            reconnect_count, reconnect_delay = 0, FIRST_RECONNECT_DELAY
            while reconnect_count < MAX_RECONNECT_COUNT:
                print("Reconnecting in {} seconds...".format(reconnect_delay))
                time.sleep(reconnect_delay)

                try:
                    client.reconnect()
                    print("Reconnected successfully!")
                    return
                except Exception as err:
                    print("{}. Reconnect failed. Retrying...".format(err))

                reconnect_delay *= RECONNECT_RATE
                reconnect_delay = min(reconnect_delay, MAX_RECONNECT_DELAY)
                reconnect_count += 1
            print("Reconnect failed after {} attempts. Exiting...".format(reconnect_count))

        def on_message(client, userdata, msg):
            """
            Callback triggered when a message is received on the subscribed topic.
            """
            print("Received `{}` from `{}` topic".format(msg.payload.decode(), msg.topic))
            subprocess.run(["./sensor_data"])  # Execute the sensor_data script upon message receipt

        # Create an MQTT client instance, assign the callback functions and connect
        client = mqtt_client.Client(client_id)
        client.on_connect = on_connect
        client.on_disconnect = on_disconnect
        client.on_message = on_message
        client.connect(broker, port)
        return client  # Return the configured client instance

    def run():
        """
        Main function to connect the MQTT client and start its loop.
        """
        client = connect_mqtt()
        client.loop_forever()

    if __name__ == '__main__':
        run()
\end{minted}

\usubsection{Sensor Simulator Data Generation and Publishing script}
\begin{minted}[%
    framesep=2mm,
    baselinestretch=1.2,
    bgcolor=LightGray,
    fontsize=\footnotesize,
    label=lst:sensor simulator data generation and publishing,
    breaklines
    ]{python}
    import json
    import random, os
    import socket
    import subprocess
    from dotenv import load_dotenv
    import paho.mqtt.client as mqtt_client

    port = 1883  # Default port for MQTT communication
    hostname = socket.gethostname()  # Retrieve the hostname of the current machine. Hostname is enforced by docker compose.
    topic = "sensor-data/{}".format(hostname)  # Define the unique topic for publishing sensor data
    client_id = 'publish-{}'.format(hostname)  # Unique client ID for MQTT connection
    # broker = 'localhost'  # Uncomment this line for local broker testing
    broker = 'host.docker.internal'  # Use this broker when running in a Docker environment

    def connect_mqtt():
        """
        Connects to the MQTT broker and sets up the on_connect callback.
        """

        def on_connect(client, userdata, flags, rc):
            """
            Callback triggered upon connecting to the MQTT broker.
            """
            if rc == 0:
                print("Connected to MQTT Broker!")
            else: 
                print("Failed to connect, return code %d\n", rc)

        client = mqtt_client.Client(client_id)
        client.on_connect = on_connect 
        client.connect(broker, port)  
        return client  

    def data_gen():
        """
        Generates sensor data by selecting a random entry from a dataset.
        The selection point is controlled via an environment variable.
        """
        with open("dataset.json") as datafile:
            json_data = json.load(datafile)  # Load the dataset from the JSON file
        
            load_dotenv()  # Load environment variables from the .env file
            startpoint = int(os.getenv('STARTPOINT'))  # Get the STARTPOINT from the .env file
            # Randomize the data selection around the startpoint. Slowly moves forward in time, while keeping results semi-random and ensuring a logical history.
            randomizer = random.randint(startpoint - 2, startpoint + 3)
            while randomizer not in range(1, len(json_data)):  # Ensure the randomizer is valid
                subprocess.run(["./set_startpoint"])  # Run a script to reset the startpoint
                load_dotenv()  # Reload environment variables
                startpoint = int(os.getenv('STARTPOINT'))
                randomizer = random.randint(startpoint - 10, startpoint + 10)

            # Update the STARTPOINT in the .env file
            with open(".env", "w") as f:
                f.write("STARTPOINT={}".format(randomizer))

            randata = json_data[str(randomizer)]  # Fetch the random data entry
            return randata  # Return the selected data

    def publish(client, data):
        """
        Publishes a message to the MQTT topic.
        """
        msg = str(data)  # Convert the data to a string
        result = client.publish(topic, msg)  # Publish the message to the topic
        status = result[0]  # Check the result status
        if status == 0:
            print("Sent `{}` to topic `{}`".format(msg, topic))
        else:
            print("Failed to send message to topic {}".format(topic))        

    if __name__ == '__main__':
        client = connect_mqtt()  # Establish the MQTT connection
        client.loop_start()  # Start the MQTT client loop in a separate thread
        randata = data_gen()  # Generate random sensor data
        publish(client, randata)  # Publish the generated data to the topic
        client.loop_stop()  # Stop the MQTT client loop
\end{minted}

\usubsection{Startpoint Setter Script}
\begin{minted}[%
    framesep=2mm,
    baselinestretch=1.2,
    bgcolor=LightGray,
    fontsize=\footnotesize,
    breaklines
    ]{python}
    import json, random

    def set_startpoint():
        """
        Sets a STARTPOINT value based on the dataset's size and writes it to an .env file.
        """
        with open("dataset.json") as datafile:
            json_data = json.load(datafile) # Load the dataset from a JSON file
            
            endpoint = round(len(json_data)/10) # Determine the upper limit for the STARTPOINT range
            startpoint = str(random.randint(1, endpoint))

            print("Setting STARTPOINT as {}".format(startpoint))

            with open(".env", "w") as f:
                f.write("STARTPOINT={}".format(startpoint)) # Write the STARTPOINT to the .env file

    if __name__ == '__main__':
        set_startpoint()
\end{minted}

\usubsection{Sensor Simulator Dockerfile}
\begin{minted}[%
    framesep=2mm,
    baselinestretch=1.2,
    bgcolor=LightGray,
    fontsize=\footnotesize,
    breaklines
    ]{python}
    import json, random

    def set_startpoint():
        """
        Sets a STARTPOINT value based on the dataset's size and writes it to an .env file.
        """
        with open("dataset.json") as datafile:
            json_data = json.load(datafile) # Load the dataset from a JSON file
            
            endpoint = round(len(json_data)/10) # Determine the upper limit for the STARTPOINT range
            startpoint = str(random.randint(1, endpoint))

            print("Setting STARTPOINT as {}".format(startpoint))

            with open(".env", "w") as f:
                f.write("STARTPOINT={}".format(startpoint)) # Write the STARTPOINT to the .env file

    if __name__ == '__main__':
        set_startpoint()
\end{minted}

\usubsection{Controller Node main application}
\begin{minted}[%
    framesep=2mm,
    baselinestretch=1.2,
    bgcolor=LightGray,
    fontsize=\footnotesize,
    breaklines
    ]{python}
    import time
    from fastapi import FastAPI, Response
    from prometheus_client import generate_latest
    from exporter_func import *

    app = FastAPI(debug=False)

    # Define an endpoint to serve metrics
    @app.get("/metrics")
    def get_metrics_app():
        start_time = time.time()
        
        # Connect to MQTT and start collecting data
        sub_client = connect_sub_mqtt()
        sub_client.loop_start()
        gather_data()
        time.sleep(1) # Allow time for data collection, in case sensors take time to respond
        sub_client.loop_stop()

        # Generate metrics data in Prometheus format
        data = generate_latest()

        end_time = time.time()
        execution_time = end_time - start_time
        print("Execution time:",execution_time)

        # Return the metrics data as a plain text HTTP response
        return Response(content=data, media_type="text/plain")

    # Define a root endpoint with a simple message
    @app.get("/")
    async def root():
        return "Go to /metrics for gitlab metrics"

    # Note for running the application locally
    # Use the command:
    # gunicorn -b localhost:8000 exporter:app -k uvicorn.workers.UvicornWorker
\end{minted}

\usubsection{Controller Node Data collection and Exporter Functions}
\begin{minted}[%
    framesep=2mm,
    baselinestretch=1.2,
    bgcolor=LightGray,
    fontsize=\footnotesize,
    breaklines
    ]{python}
    import paho.mqtt.client as mqtt_client
    import json, time
    from exporter_var import *

    # Function to establish a connection to the MQTT broker and subscribe to a topic
    def connect_sub_mqtt():

        # Callback for successful connection to the MQTT broker
        def on_connect(client, userdata, flags, rc):
            if rc == 0:
                print("Connected to MQTT Broker!")
                client.subscribe(sub_topic) # Subscribe to the specified topic
            else:
                print("Failed to connect, return code %d\n", rc)
        
        # Callback for handling disconnection from the MQTT broker
        def on_disconnect(client, userdata, rc):

            FIRST_RECONNECT_DELAY = 1
            RECONNECT_RATE = 2
            MAX_RECONNECT_COUNT = 12
            MAX_RECONNECT_DELAY = 60

            print("Disconnected with result code: %s", rc)
            reconnect_count, reconnect_delay = 0, FIRST_RECONNECT_DELAY
            while reconnect_count < MAX_RECONNECT_COUNT:
                print("Reconnecting in {} seconds...".format(reconnect_delay))
                time.sleep(reconnect_delay)

                try:
                    client.reconnect()
                    print("Reconnected successfully!")
                    return
                except Exception as err:
                    print("{}. Reconnect failed. Retrying...".format(err))

                reconnect_delay *= RECONNECT_RATE # Increase delay exponentially
                reconnect_delay = min(reconnect_delay, MAX_RECONNECT_DELAY)
                reconnect_count += 1
            print("Reconnect failed after {} attempts. Exiting...".format(reconnect_count))

        # Callback for receiving messages from the MQTT broker
        def on_message(client, userdata, msg):

            print("Received data from `{}` topic".format(msg.topic))
            promethify_data(msg)  # Process the received message for Prometheus metrics

        # Initialize the MQTT client and assign the callbacks
        client = mqtt_client.Client(sub_client_id)
        client.on_connect = on_connect
        client.on_disconnect = on_disconnect
        client.on_message = on_message
        client.connect(broker, sub_port)
        return client

    # Function to continuously run the MQTT subscriber
    def sub_client_run():
        sub_client = connect_sub_mqtt()
        sub_client.loop_forever()

    # Function to publish a request message to collect data
    def request_data(client):
        msg = 'Time to collect data.'
        result = client.publish(pub_topic, msg)
        status = result[0]
        if status == 0:
            print("Sent `{}` to topic `{}`".format(msg, pub_topic))
        else:
            print("Failed to send message to topic {}".format(pub_topic))

    # Function to connect to the MQTT broker as a publisher
    def connect_pub_mqtt():

        def on_connect(client, userdata, flags, rc):
            if rc == 0:
                print("Connected to MQTT Broker!")
            else:
                print("Failed to connect, return code %d\n", rc)

        client = mqtt_client.Client(pub_client_id)
        client.on_connect = on_connect
        client.connect(broker, pub_port)
        return client

    # Function to gather data by publishing a request message
    def gather_data():
        pub_client = connect_pub_mqtt()
        pub_client.loop_start()
        request_data(pub_client)
        pub_client.loop_stop()

    # Function to assign received data to Prometheus metrics
    def assign_data(prom_var, data_name, sensor, jsondata):

        value =  float(jsondata[data_name].replace(',', '.')) # Parse the value
        if  value != -200:  # Exclude no readings
            prom_var.labels(sensor).set(value) # Set the metric value with sensor label  

    # Function to process received MQTT messages and update Prometheus metrics
    def promethify_data(msg):

        sensor = msg.topic.split('/')[1] # Extract sensor name from topic
        # Decode and parse the JSON payload
        jsondata_string = msg.payload.decode().replace("'", '"') 
        jsondata = json.loads(jsondata_string)

        # Assign each metric to the corresponding Prometheus variable
        assign_data(air_quality_co_gt_gauge, "CO(GT)", sensor, jsondata)
        assign_data(air_quality_pt08s1_co_gauge, "PT08.S1(CO)", sensor, jsondata)
        assign_data(air_quality_nmhc_gt_gauge, "NMHC(GT)", sensor, jsondata)
        assign_data(air_quality_c6h6_gt_gauge, "C6H6(GT)", sensor, jsondata)
        assign_data(air_quality_pt08s2_nmhc_gauge, "PT08.S2(NMHC)", sensor, jsondata)
        assign_data(air_quality_nox_gt_gauge, "NOx(GT)", sensor, jsondata)
        assign_data(air_quality_pt08s3_nox_gauge, "PT08.S3(NOx)", sensor, jsondata)
        assign_data(air_quality_no2_gt_gauge, "NO2(GT)", sensor, jsondata)
        assign_data(air_quality_pt08s4_no2_gauge, "PT08.S4(NO2)", sensor, jsondata)
        assign_data(air_quality_pt08s5_o3_gauge, "PT08.S5(O3)", sensor, jsondata)
        assign_data(air_quality_t_gauge, "T", sensor, jsondata)
        assign_data(air_quality_rh_gauge, "RH", sensor, jsondata)
        assign_data(air_quality_ah_gauge, "AH", sensor, jsondata)
\end{minted}

\usubsection{Controller Node Prometheus Metrics Variables}
\begin{minted}[%
    framesep=2mm,
    baselinestretch=1.2,
    bgcolor=LightGray,
    fontsize=\footnotesize,
    breaklines
    ]{python}
    from prometheus_client import Gauge

    sub_port = 1883 # Port for the subscriber
    pub_port = 1883 # Port for the publisher
    sub_topic = "sensor-data/#" # Subscription topic for sensor data (wildcard for all subtopics)
    pub_topic = "collect-data" # Topic to publish data collection requests
    sub_client_id = 'subscribe-exporter' # Client ID for the MQTT subscriber
    pub_client_id = 'publish-exporter' # Client ID for the MQTT publisher
    #broker = 'localhost' # uncomment for local testing
    broker = 'host.docker.internal' # Docker-specific hostname for connecting to the broker


    # Prometheus metrics definitions
    # Each metric is defined as a Prometheus Gauge with a description and a label for 'sensor'

    air_quality_co_gt_gauge = Gauge('air_quality_co_gt_gauge', 'True hourly averaged concentration CO in mg/m^3 (reference analyzer)', ['sensor'])

    air_quality_pt08s1_co_gauge = Gauge('air_quality_pt08s1_co_gauge', 'Tin oxide hourly averaged sensor response (nominally CO targeted)', ['sensor'])

    air_quality_nmhc_gt_gauge = Gauge('air_quality_nmhc_gt_gauge', 'True hourly averaged overall Non Metanic HydroCarbons concentration in microg/m^3 (reference analyzer)', ['sensor'])

    air_quality_c6h6_gt_gauge = Gauge('air_quality_c6h6_gt_gauge', 'True hourly averaged Benzene concentration in microg/m^3 (reference analyzer)', ['sensor'])

    air_quality_pt08s2_nmhc_gauge = Gauge('air_quality_pt08s2_nmhc_gauge', 'Titania hourly averaged sensor response (nominally NMHC targeted)', ['sensor'])

    air_quality_nox_gt_gauge = Gauge('air_quality_nox_gt_gauge', 'True hourly averaged NOx concentration in ppb (reference analyzer)', ['sensor'])

    air_quality_pt08s3_nox_gauge = Gauge('air_quality_pt08s3_nox_gauge', 'Tungsten oxide hourly averaged sensor response (nominally NOx targeted)', ['sensor'])

    air_quality_no2_gt_gauge = Gauge('air_quality_no2_gt_gauge', 'True hourly averaged NO2 concentration in microg/m^3 (reference analyzer)', ['sensor'])

    air_quality_pt08s4_no2_gauge = Gauge('air_quality_pt08s4_no2_gauge', 'Tungsten oxide hourly averaged sensor response (nominally NO2 targeted)', ['sensor'])

    air_quality_pt08s5_o3_gauge = Gauge('air_quality_pt08s5_o3_gauge', 'Indium oxide hourly averaged sensor response (nominally O3 targeted)', ['sensor'])

    air_quality_t_gauge = Gauge('air_quality_t_gauge', 'Temperature in C', ['sensor'])

    air_quality_rh_gauge = Gauge('air_quality_rh_gauge', 'Relative Humidity (%)', ['sensor'])

    air_quality_ah_gauge = Gauge('air_quality_ah_gauge', 'Absolute Humidity', ['sensor'])
\end{minted}

\usubsection{Controller Node Dockerfile}
\begin{minted}[%
    framesep=2mm,
    baselinestretch=1.2,
    bgcolor=LightGray,
    fontsize=\footnotesize,
    breaklines
    ]{dockerfile}
    FROM python:3.11.0

    # Set the working directory inside the container
    WORKDIR /srv/exporter

    # Add the main application and supporting scripts to the working directory
    ADD exporter.py /srv/exporter/
    ADD exporter_func.py /srv/exporter/
    ADD exporter_var.py /srv/exporter/
    ADD requirements.txt /srv/exporter/

    # Set execute permissions for all users on the project directory
    RUN chmod -R a+x /srv/exporter

    # Install the Python dependencies specified in requirements.txt
    RUN python3 -m pip install -r requirements.txt

    # Expose the necessary ports for the Prometheus metrics endpoint and the MQTT protocol
    EXPOSE 8000
    EXPOSE 1883

    # Set the entrypoint command to run the application using Gunicorn with Uvicorn workers
    ENTRYPOINT [ "gunicorn", "-b", "0.0.0.0:8000", "exporter:app", "-k", "uvicorn.workers.UvicornWorker" ]
\end{minted}

\usubsection{Prometheus Configuration File}
\begin{minted}[%
    framesep=2mm,
    baselinestretch=1.2,
    bgcolor=LightGray,
    fontsize=\footnotesize,
    breaklines
    ]{yaml}
    global:
    scrape_interval:     15s
    evaluation_interval: 15s

    rule_files:
    - prometheus_alerts_rules.yml

    scrape_configs:

    - job_name: prometheus # Collecting data from Prometheus itself
        static_configs:
        - targets: ['127.0.0.1:9090']

    - job_name: sensors # Collecting metrics from the Controller Node
        scrape_interval: 60s
        scrape_timeout: 50s
        static_configs:
        - targets: ['host.docker.internal:8000'] # Using host.docker.internal to direct containerized prometheus to host's 127.0.0.1 (localhost)
\end{minted}

\usubsection{Docker Compose File}
\begin{minted}[%
    framesep=2mm,
    baselinestretch=1.2,
    bgcolor=LightGray,
    fontsize=\footnotesize,
    breaklines
    ]{yaml}
services:
    mqtt_broker:
        image: "eclipse-mosquitto:2.0.18"  # Using the Eclipse Mosquitto MQTT broker
        container_name: "thesis-app.mqtt_broker"
        volumes: 
        - ./mqtt_broker/mosquitto.conf:/mosquitto/config/mosquitto.conf  # Mount custom Mosquitto configuration
        ports:
        - "1883:1883"  # MQTT default port
        - "9001:9001"  # WebSockets support for MQTT
        restart: unless-stopped  # Automatically restart unless manually stopped

    exporter:
        image: "anastzampetis/iot-exporter"  # Custom IoT data exporter
        container_name: "thesis-app.exporter"
        ports:
        - "8000:8000"  # Expose exporter service on port 8000
        restart: unless-stopped  # Ensure the service restarts if it fails

    prometheus:
        image: "prom/prometheus:v2.46.0"  # Prometheus monitoring service
        container_name: "thesis-app.prometheus"
        volumes:
        - ./prometheus/prometheus.yml:/etc/prometheus/prometheus.yml  # Custom Prometheus config
        - ./prometheus/prometheus_alerts_rules.yml:/etc/prometheus/prometheus_alerts_rules.yml  # Alert rules (Empty but needed as file)
        - prometheus_data:/prometheus  # Persistent data storage
        ports:
        - "9090:9090"  # Expose Prometheus on port 9090
        command:
        - "--config.file=/etc/prometheus/prometheus.yml"  # Specify config file
        - "--web.enable-lifecycle"  # Allow configuration reload via API
        - "--storage.tsdb.path=/prometheus"  # Set data storage location
        - "--storage.tsdb.retention.time=1y"  # Retention policy. Keep data for 1 year
        restart: unless-stopped  

    grafana:
        image: "grafana/grafana:10.2.2"  # Grafana visualization service
        container_name: "thesis-app.grafana"
        ports:
        - "3000:3000"  # Expose Grafana GUI on port 3000
        restart: unless-stopped 
        environment:
        - GF_SECURITY_ADMIN_USER=admin  # Admin username for Grafana
        - GF_SECURITY_ADMIN_PASSWORD=grafana  # Admin password for Grafana
        - GF_PATHS_DATA=/var/lib/grafana  # Set data storage path
        volumes:
        - ./grafana:/var/lib/grafana  # Persist Grafana data

    # Sensor services - each sensor must be initialized separately because Docker replicas do not support dynamic variables and container names are not accessible from inside the container. (needed for MQQT client)
    sensor-1:
        image: "anastzampetis/sensor-emul"  # Custom sensor emulator image
        container_name: "thesis-app.sensor-1"
        hostname: "sensor-1"  # Set hostname for usage in the MQTT client
        restart: unless-stopped 

    sensor-2:
        image: "anastzampetis/sensor-emul"
        container_name: "thesis-app.sensor-2"
        hostname: "sensor-2"
        restart: unless-stopped

    sensor-3:
        image: "anastzampetis/sensor-emul"
        container_name: "thesis-app.sensor-3"
        hostname: "sensor-3"
        restart: unless-stopped

    sensor-4:
        image: "anastzampetis/sensor-emul"
        container_name: "thesis-app.sensor-4"
        hostname: "sensor-4"
        restart: unless-stopped

    sensor-5:
        image: "anastzampetis/sensor-emul"
        container_name: "thesis-app.sensor-5"
        hostname: "sensor-5"
        restart: unless-stopped

volumes:
    prometheus_data:  # Persistent volume for Prometheus data storage
\end{minted}